\chapter{Flat Feet (Pes Planus)}

\section*{Definition}
Pes planus, or flat feet, is a postural deformity characterized by collapse of the medial longitudinal arch of the foot, resulting in the entire sole contacting the ground. It may be flexible (arch reappears with non-weight-bearing) or rigid, and is symptomatic when associated with pain, fatigue, or functional impairment.

\section*{What Happens in Your Body}
The medial longitudinal arch is maintained by the bony architecture of the talus, navicular, and cuneiforms, along with the spring (calcaneonavicular) ligament, plantar fascia, and the posterior tibial tendon. Dysfunction of the posterior tibial tendon (the primary dynamic supporter) leads to progressive arch collapse, forefoot abduction, and hindfoot valgus. Chronic strain on the plantar fascia and ligamentous structures produces pain and secondary osteoarthritis.

\section*{Causes}
\begin{itemize}
  \item \textbf{Congenital:} Hypermobility, ligamentous laxity, tarsal coalition, accessory navicular
  \item \textbf{Idiopathic:} Flexible flatfoot in children (often resolves with growth)
  \item \textbf{Adult-acquired:} Posterior tibial tendon dysfunction (PTTD), often in middle-aged women
  \item \textbf{Traumatic:} Spring ligament tear, tarsometatarsal (Lisfranc) injury, calcaneal fracture malunion
  \item \textbf{Neuromuscular:} Cerebral palsy, muscular dystrophy, Charcot-Marie-Tooth disease
  \item \textbf{Inflammatory:} Rheumatoid arthritis, seronegative spondyloarthropathy
  \item \textbf{Systemic:} Obesity, diabetes, Ehlers-Danlos syndrome, pregnancy-related ligamentous laxity
\end{itemize}

\section*{When to See a Doctor}
See your doctor if:
\begin{itemize}
  \item Persistent foot pain, fatigue, or swelling along the medial arch
  \item Difficulty walking or standing for prolonged periods
  \item Asymmetric flattening or sudden change in foot shape
  \item Inability to perform a single-leg heel rise (suggests PTTD)
\end{itemize}

\section*{Related Symptoms}
\begin{itemize}
  \item Medial arch and ankle pain
  \item Hindfoot valgus (too-many-toes sign)
  \item Forefoot abduction (lateral deviation of the forefoot)
  \item Plantar fasciitis and Achilles tendonitis
  \item Sinus tarsi syndrome (lateral foot pain)
\end{itemize}
\textbf{Important:} Symptomatic adult-acquired flatfoot is most commonly due to posterior tibial tendon dysfunction and should be staged radiographically to guide management.

\section*{Self-Care / Home Management}
\begin{itemize}
  \item Supportive footwear with arch support and firm heel counter
  \item Over-the-counter orthotic inserts (arch supports, heel cups)
  \item Activity modification to reduce prolonged standing or walking on hard surfaces
  \item Ice massage to the medial arch after activity
  \item Calf stretching and intrinsic foot muscle strengthening (towel curls, marble pickups)
\end{itemize}

\section*{How Your Doctor Will Evaluate You}
\begin{itemize}
  \item Gait observation: heel rise test, too-many-toes sign, hindfoot alignment
  \item Weight-bearing radiography: lateral talo-metatarsal angle, Meary angle, calcaneal pitch
  \item MRI to evaluate posterior tibial tendon integrity and spring ligament
  \item Ultrasound for dynamic assessment of tendon function
  \item CT for evaluation of tarsal coalition or degenerative changes
\end{itemize}

\section*{Treatment Options}
\begin{itemize}
  \item Activity modification, physical therapy, and orthotic management for flexible flatfoot
  \item NSAIDs for pain and inflammation in acute flares
  \item Corticosteroid injection around the posterior tibial tendon (avoid intratendinous injection)
  \item Ankle-foot orthosis or Arizona brace for advanced PTTD
  \item Surgical reconstruction: tendon transfer, calcaneal osteotomy, arthrodesis for rigid or refractory deformity
\end{itemize}

\section*{References}
\begin{thebibliography}{99}
\bibitem{ref1} Harris EJ, Vanore JV, Thomas JL, et al. "Diagnosis and Treatment of Pediatric Flatfoot." J Foot Ankle Surg. 2004;43(6):341--373.
\bibitem{ref2} Mosca VS. "Flexible Flatfoot in Children and Adolescents." J Child Orthop. 2010;4(2):107--121.
\bibitem{ref3} Sullivan JA. "Pediatric Flatfoot: Evaluation and Management." J Am Acad Orthop Surg. 1999;7(1):44--53.
\bibitem{ref4} Evans AM, Rome K. "A Cochrane Review of the Evidence for Nonsurgical Interventions for Pediatric Flatfoot." Pediatr Phys Ther. 2012;24(1):91--95.
\bibitem{ref5} Toullec E. "Adult Flatfoot." Orthop Traumatol Surg Res. 2015;101(1 Suppl):S73--S81.
\end{thebibliography}
\clearpage
